\section{Modelo de memória}
\label{sec:modelo-memoria}

O capítulo~17 da JLS define o que uma thread pode ver das escritas de
outra. O conceito central é a relação
\termo{happens-before}{\enquote{Acontece antes}: se a escrita A
happens-before a leitura B, B enxerga o efeito de A. Sem essa relação
entre duas ações em threads diferentes, não há garantia de visibilidade
nem de ordem.} (§17.4.5). Exemplos da própria JLS: a escrita num campo
\texttt{volatile} happens-before toda leitura subsequente desse campo, e
um \texttt{unlock} de um monitor happens-before todo \texttt{lock}
subsequente do mesmo monitor. Campos \texttt{final} têm regras próprias
de publicação segura (§17.5).

Se um programa não tem \termo{data race}{Duas threads acessam a mesma
variável, pelo menos uma escreve, e não há happens-before entre os
acessos.}, ele se comporta como se fosse sequencial; com data races, uma
leitura pode ver valores surpreendentes. O texto de Shipilëv
(\citetitle{artigo:shipilev-jmm}) desenvolve essas consequências.

\paragraph{O bug que o bytecode mostra.}
Na listagem de \texttt{Sincronizacao} (seção~\ref{sec:sincronizacao}),
\texttt{contador++} ocupa os offsets 6 a 11: \texttt{getfield} lê o
campo, \texttt{iconst\_1} e \texttt{iadd} somam na pilha, e
\texttt{putfield} grava. São três passos, não um.
Entre o \texttt{getfield} e o \texttt{putfield}, outra thread pode ler o
mesmo valor antigo; as duas gravam o mesmo resultado e um incremento se
perde. \texttt{volatile} não resolve: ele garante visibilidade, mas o
ler-somar-gravar continua sendo três passos. As opções corretas são
\texttt{synchronized} (o que o exemplo faz), \texttt{AtomicInteger} ou
\texttt{LongAdder}, ou não compartilhar o estado. É o tema do item 78
do \citetitle{livro:bloch2017} (\enquote{Synchronize access to shared
mutable data}).

\paragraph{A decisão que isso muda.}
Records e campos \texttt{final} não são só estilo: objetos imutáveis
podem ser compartilhados entre threads sem lock (JLS §17.5 dá a garantia
para os campos \texttt{final}). Num serviço web, a pergunta \enquote{este
estado é compartilhado entre requisições?} vem antes de \enquote{qual
lock usar?}. O \citetitle{livro:goetz2006} é a referência para o resto.
